\begin{lemma}[The complex Riesz transform is $\mathbb{C}$-linear, and unimodular rotations do not change the pointwise size of its values]
Let $R_\gamma^{\mathbb{C}} : L^2(\gamma_n;\mathbb{C}) \to L^2(\gamma_n;\mathbb{C}^n)$
be the complex Riesz transform of \noderef{ComplexRieszTransform}.  Then:
\begin{enumerate}
\item (additivity) $R_\gamma^{\mathbb{C}}(v + w)
  = R_\gamma^{\mathbb{C}} v + R_\gamma^{\mathbb{C}} w$ for all
  $v, w \in L^2(\gamma_n;\mathbb{C})$;
\item ($\mathbb{C}$-homogeneity) $R_\gamma^{\mathbb{C}}(c\,v)
  = c\,R_\gamma^{\mathbb{C}} v$ for every $c \in \mathbb{C}$ and every
  $v \in L^2(\gamma_n;\mathbb{C})$;
\item (unimodular invariance of the pointwise size) if $c \in \mathbb{C}$ with
  $|c| = 1$, then for every $v \in L^2(\gamma_n;\mathbb{C})$ and
  $\gamma_n$-almost every $x \in \mathbb{R}^n$,
\[
  \bigl|R_\gamma^{\mathbb{C}}(c\,v)(x)\bigr|_{\mathbb{C}^n}
    = \bigl|R_\gamma^{\mathbb{C}} v(x)\bigr|_{\mathbb{C}^n} .
\]
\end{enumerate}
Clauses (1) and (2) say that $R_\gamma^{\mathbb{C}}$ is the complex-linear
extension of $R_\gamma$ that the reference paper's complex-valued estimate is
about, and clause (3) is the sentence ``Complex linearity implies
$|R_\gamma(e^{i\theta}f)| = |R_\gamma f|$'' that licenses the paper's rotation
average; it is stated for an arbitrary unimodular $c$, of which $c = e^{i\theta}$
is the case the paper uses.  Clause (3) is recorded here, rather than left to
the reader of clause (2), because it is an assertion about the values of
representatives at almost every point, whereas clause (2) is an identity between
equivalence classes.

The scalars in clause (2) are the \emph{complex} ones.  This is strictly more
than the real linearity of $R_\gamma$ in \noderef{RieszTransformLinear}: the
real and imaginary parts of $c\,v$ are not $c$ times those of $v$ --- already for
$c = i$ one has $\operatorname{Re}(iv) = -\operatorname{Im} v$ and
$\operatorname{Im}(iv) = \operatorname{Re} v$, so the two parts are exchanged ---
and clause (2) is therefore not a restatement of real homogeneity but a genuine
consequence of it combined with the structure of
\noderef{ComplexRieszTransform}.
\end{lemma}
\begin{proof}
Throughout write $\iota : \mathbb{R}^n \to \mathbb{C}^n$ for the coordinatewise
inclusion of \noderef{ComplexRieszTransform} and recall from that node the
defining formula
\begin{equation*}
  R_\gamma^{\mathbb{C}} z
    = \iota\bigl(R_\gamma(\operatorname{Re} z)\bigr)
      + i\,\iota\bigl(R_\gamma(\operatorname{Im} z)\bigr),
  \qquad z \in L^2(\gamma_n;\mathbb{C}) . \tag{D}
\end{equation*}
Three elementary facts are used repeatedly.  First, the maps
$z \mapsto \operatorname{Re} z$ and $z \mapsto \operatorname{Im} z$ from
$L^2(\gamma_n;\mathbb{C})$ to $L^2(\gamma_n;\mathbb{R})$ are additive and
$\mathbb{R}$-homogeneous, because they are composition with the
$\mathbb{R}$-linear maps $\operatorname{Re}, \operatorname{Im} : \mathbb{C} \to
\mathbb{R}$ and the vector operations on a Lebesgue space are performed on
representatives.  Second, the map $y \mapsto \iota(y)$ from
$L^2(\gamma_n;\mathbb{R}^n)$ to $L^2(\gamma_n;\mathbb{C}^n)$ is additive and
$\mathbb{R}$-homogeneous, for the same reason, $\iota$ being
$\mathbb{R}$-linear.  Third, $L^2(\gamma_n;\mathbb{C}^n)$ is a complex vector
space whose real scalar action is the restriction of its complex one, so real
and complex scalars may be combined freely.

\emph{Step 1: clause (1).}
Let $v, w \in L^2(\gamma_n;\mathbb{C})$.  By the first fact,
$\operatorname{Re}(v+w) = \operatorname{Re} v + \operatorname{Re} w$ and
$\operatorname{Im}(v+w) = \operatorname{Im} v + \operatorname{Im} w$.  By the
additivity clause of \noderef{RieszTransformLinear}, applied to the two real
inputs in each line,
\[
  R_\gamma\bigl(\operatorname{Re}(v+w)\bigr)
    = R_\gamma(\operatorname{Re} v) + R_\gamma(\operatorname{Re} w),
  \qquad
  R_\gamma\bigl(\operatorname{Im}(v+w)\bigr)
    = R_\gamma(\operatorname{Im} v) + R_\gamma(\operatorname{Im} w).
\]
Substituting into (D) at $z = v + w$ and using the second and third facts to
distribute $\iota$ and the factor $i$ over the sums,
\[
  R_\gamma^{\mathbb{C}}(v+w)
    = \Bigl(\iota R_\gamma(\operatorname{Re} v)
        + i\,\iota R_\gamma(\operatorname{Im} v)\Bigr)
      + \Bigl(\iota R_\gamma(\operatorname{Re} w)
        + i\,\iota R_\gamma(\operatorname{Im} w)\Bigr)
    = R_\gamma^{\mathbb{C}} v + R_\gamma^{\mathbb{C}} w ,
\]
the regrouping of the four summands being commutativity and associativity of
addition in $L^2(\gamma_n;\mathbb{C}^n)$, and the two bracketed expressions
being $R_\gamma^{\mathbb{C}} v$ and $R_\gamma^{\mathbb{C}} w$ by (D).

\emph{Step 2: the parts of $c\,v$.}
Let $c \in \mathbb{C}$ and put $a = \operatorname{Re} c$, $b = \operatorname{Im} c$,
so $c = a + ib$ with $a, b$ real.  For any $z \in \mathbb{C}$,
\[
  c\,z = (a + ib)(\operatorname{Re} z + i\operatorname{Im} z)
    = \bigl(a\operatorname{Re} z - b\operatorname{Im} z\bigr)
      + i\bigl(a\operatorname{Im} z + b\operatorname{Re} z\bigr),
\]
and the two bracketed quantities are real, so
$\operatorname{Re}(cz) = a\operatorname{Re} z - b\operatorname{Im} z$ and
$\operatorname{Im}(cz) = a\operatorname{Im} z + b\operatorname{Re} z$.  Applying
this at $z = v(x)$ for almost every $x$, and using that the vector operations on
$L^2$ are representativewise, we get the identities in
$L^2(\gamma_n;\mathbb{R})$
\[
  \operatorname{Re}(c\,v) = a\,\operatorname{Re} v - b\,\operatorname{Im} v,
  \qquad
  \operatorname{Im}(c\,v) = a\,\operatorname{Im} v + b\,\operatorname{Re} v .
\]

\emph{Step 3: transforming the parts.}
Write $u_1 = R_\gamma(\operatorname{Re} v)$ and
$u_2 = R_\gamma(\operatorname{Im} v)$.  Both clauses of
\noderef{RieszTransformLinear} are unconditional, so they may be applied to the
combinations of Step 2 with the real scalars $a$, $-b$ and $b$ (additivity for
the sum, homogeneity for each scalar multiple, a difference being the sum with
the scalar multiple by $-1$):
\[
  R_\gamma\bigl(\operatorname{Re}(c\,v)\bigr) = a\,u_1 - b\,u_2,
  \qquad
  R_\gamma\bigl(\operatorname{Im}(c\,v)\bigr) = a\,u_2 + b\,u_1 .
\]

\emph{Step 4: clause (2).}
Substituting Step 3 into (D) at $z = c\,v$ and distributing $\iota$ and the
factor $i$ by the second and third facts,
\[
  R_\gamma^{\mathbb{C}}(c\,v)
    = \iota\bigl(a\,u_1 - b\,u_2\bigr) + i\,\iota\bigl(a\,u_2 + b\,u_1\bigr)
    = a\,\iota u_1 - b\,\iota u_2 + ia\,\iota u_2 + ib\,\iota u_1 .
\]
Collecting the two terms carrying $\iota u_1$ and the two carrying $\iota u_2$,
\[
  R_\gamma^{\mathbb{C}}(c\,v)
    = (a + ib)\,\iota u_1 + (ia - b)\,\iota u_2 .
\]
Now $ia - b = i(a + ib)$, because $i(a+ib) = ia + i^2 b = ia - b$.  Hence
\[
  R_\gamma^{\mathbb{C}}(c\,v)
    = c\,\iota u_1 + i\,c\,\iota u_2
    = c\,\bigl(\iota u_1 + i\,\iota u_2\bigr)
    = c\,R_\gamma^{\mathbb{C}} v ,
\]
the last equality by (D) at $z = v$.  This is clause (2).  Taking $c = i$, that
is $a = 0$ and $b = 1$, the computation reads
$R_\gamma^{\mathbb{C}}(iv) = -\iota u_2 + i\,\iota u_1 = i(\iota u_1 + i\iota u_2)
= i R_\gamma^{\mathbb{C}} v$, and it is visibly the real homogeneity of
$R_\gamma$ --- through $R_\gamma(-\operatorname{Im} v) = -u_2$ --- that makes it
work.

\emph{Step 5: clause (3).}
Let $|c| = 1$ and $v \in L^2(\gamma_n;\mathbb{C})$.  By clause (2),
$R_\gamma^{\mathbb{C}}(c\,v) = c\,R_\gamma^{\mathbb{C}} v$ as elements of
$L^2(\gamma_n;\mathbb{C}^n)$, that is, as $\gamma_n$-a.e.\ equivalence classes.
Scalar multiplication there is performed on representatives, so for
$\gamma_n$-almost every $x$,
\[
  R_\gamma^{\mathbb{C}}(c\,v)(x) = c\,\bigl(R_\gamma^{\mathbb{C}} v(x)\bigr)
  \qquad\text{in } \mathbb{C}^n .
\]
Taking norms at such an $x$ and using absolute homogeneity of the norm of
$\mathbb{C}^n$ over the complex scalars,
\[
  \bigl|R_\gamma^{\mathbb{C}}(c\,v)(x)\bigr|_{\mathbb{C}^n}
    = |c|\;\bigl|R_\gamma^{\mathbb{C}} v(x)\bigr|_{\mathbb{C}^n}
    = \bigl|R_\gamma^{\mathbb{C}} v(x)\bigr|_{\mathbb{C}^n} ,
\]
since $|c| = 1$.  This is clause (3).
\end{proof}
